\clearpage
\section*{Reproducibility and sources}
\textbf{Supplied computational files.}
The folder \texttt{experiments/} contains the serial C++17 heat solver
\texttt{heat\_serial.cpp} and the Python script \texttt{run\_experiments.py}.
The script generates the numerical tables and figures for Tasks 3.2, 3.5, 4.1, 6.7, and 7.2.
The directory \texttt{results/} contains CSV tables, raw heat-solver output, and binary snapshot fields.
The directories \texttt{tables/} and \texttt{figures/} contain the files included by this document.

On a system with a C++17 compiler, Python, NumPy, and Matplotlib, run from the document folder:
\begin{lstlisting}
g++ -std=c++17 -O3 -ffp-contract=off experiments/heat_serial.cpp -o experiments/heat_serial
python experiments/run_experiments.py
pdflatex -interaction=nonstopmode -halt-on-error Methode_Numerique_Tasks_Solutions.tex
\end{lstlisting}
On Windows, use \texttt{experiments/heat\_serial.exe} as the executable output name; the Python
script selects the corresponding executable automatically. For a compiler with a different
output location, pass its path as the script's first argument.

All reported numerical experiments use double precision. Timings recorded in the raw output
are incidental and are not used as project cluster-performance results. The seven-operation
counts in Task 6.6 are theoretical estimates under the stated sustained-rate assumption.
The instability experiment is deliberately isolated from the ordinary stability checks.

\textbf{Primary sources.}
\begin{enumerate}
\item \emph{Numerical Methods for the Heat Equation}, supplied course background notes,
Chapters 1--8. Task numbers, notation, input data, and requested experiments follow this document.
\item \emph{Parallel 2D Heat Diffusion with MPI over Tailscale}, supplied project specification,
Section 3 and Part 1. The boundary-heated initial data used in Task 6.7 and the serial-reference
implementation conventions follow this specification.
\end{enumerate}

\textbf{Assistance acknowledgement.}
OpenAI Codex assisted with the mathematical exposition, experimental C++/Python programs,
and LaTeX preparation.
